% Part: computability
% Chapter: recursive-functions
% Section: other-recursions

\documentclass[../../../include/open-logic-section]{subfiles}

\begin{document}

\olfileid{cmp}{rec}{ore}
\olsection{પુનરાવર્તનનાં અન્ય સ્વરૂપો}

જોડ-નિરૂપણ અને અનુક્રમ-નિરૂપણ વડે આદિમ પુનરાવર્તનનાં વધુ અસામાન્ય
(અને ઉપયોગી) સ્વરૂપોનું સમર્થન કરી શકીએ છીએ. દાખલા તરીકે, નીચેની
વ્યાખ્યામાં જેમ થાય છે તેમ, બે વિધેયો એકસાથે વ્યાખ્યાયિત કરવાં ઘણી
વાર ઉપયોગી છે:
\begin{align*}
h_0(\vec x, 0) & = f_0(\vec x) \\
h_1(\vec x, 0) & = f_1(\vec x) \\
h_0(\vec x, y+1) & = g_0(\vec x, y, h_0(\vec x, y), h_1(\vec x, y)) \\
h_1(\vec x, y+1) & = g_1(\vec x, y, h_0(\vec x, y), h_1(\vec x, y))
\end{align*}
આ \emph{એકસાથે પુનરાવર્તન}નું ઉદાહરણ છે. વિધેયો વ્યાખ્યાયિત
કરવાની બીજી ઉપયોગી રીતમાં, નીચે મુજબ $h(\vec x, y+1)$નું મૂલ્ય
અગાઉનાં \emph{બધાં} મૂલ્યો $h(\vec x, 0)$, \dots,~$h(\vec x, y)$ના
આધારે આપીએ છીએ:
\begin{align*}
h(\vec x, 0) & = f(\vec x) \\
h(\vec x, y+1) & = g(\vec x, y, \tuple{h(\vec x, 0), \dots, h(\vec x, y)}).
\end{align*}
આ વિચારને નીચેની યોજના વધુ ટૂંકમાં વ્યક્ત કરે છે:
\[
h(\vec x, y) = g(\vec x, y, \tuple{h(\vec x, 0), \dots, h(\vec x, y-1)})
\]
અહીં $y$ એ $0$ હોય ત્યારે $g$નું છેલ્લું આર્ગ્યુમેન્ટ માત્ર ખાલી
અનુક્રમ છે એમ સમજવું. બંને રજૂઆતોમાં વિચાર એ છે કે ``ઉત્તરગામી
તબક્કો'' ગણતી વખતે $h$ અત્યાર સુધી ગણાયેલા મૂલ્યોનો આખો અનુક્રમ
વાપરી શકે છે. તેને \emph{અગાઉનાં બધાં મૂલ્યો પર આધારિત પુનરાવર્તન}
કહે છે. ખાસ ઉદાહરણ તરીકે, તે નીચેના પ્રકારની વ્યાખ્યાનું સમર્થન કરે
છે:
\begin{align*}
h(\vec x, y) & = \begin{cases}
  g(\vec x, y, h(\vec x, k(\vec x, y))) & \text{જો $k(\vec x, y) < y$ હોય} \\
  f(\vec x) & \text{અન્યથા}
\end{cases}
\end{align*}
બીજા શબ્દોમાં, $y$ પર $h$નું મૂલ્ય~$k$ વડે અપાયેલા પહેલાંના
\emph{કોઈ પણ} સ્થાને $h$ના મૂલ્ય પર આધારિત ગણી શકાય છે.


\begin{prob}
  અગાઉનાં બધાં મૂલ્યો પર આધારિત પુનરાવર્તન વડે શેષ વિધેય
  $r(x,y)$ વ્યાખ્યાયિત કરો. ($x$, $y$ પ્રાકૃતિક સંખ્યાઓ અને
  $y > 0$ હોય, તો $r(x,y)$ એ~$y$થી નાની એવી સંખ્યા છે કે કોઈ~$z$
  માટે $x = z\times y + r(x,y)$ થાય છે. ચોક્કસતા માટે $y=0$ હોય
  ત્યારે $r(x,0) = 0$ માનીએ.)
\end{prob}

આ વિધેયો સામાન્ય આદિમ પુનરાવર્તનથી કેવી રીતે મેળવી શકાય તે
વિચારો. આદિમ પુનરાવર્તનનું એક છેલ્લું સ્વરૂપ વધુ લવચીક છે: તેમાં
ચાલતાં ચાલતાં \emph{પ્રાચલો} (બાજુનાં મૂલ્યો) બદલવાની છૂટ છે:
\begin{align*}
h(\vec x, 0) & = f(\vec x) \\
h(\vec x, y+1) & = g(\vec x, y, h(k(\vec x), y))
\end{align*}
તેને પણ સામાન્ય આદિમ પુનરાવર્તનથી અનુકરિત કરી શકાય છે. (આ કરવું
સરળ નથી. સંકેત તરીકે ગણતરીને હાથથી ઊલટી દિશામાં ઉકેલી જુઓ.)

\end{document}
